\section{ind-对象与 pro-对象}\etiquette[8]{I.sec8}\pageoriginale

\subsection{共尾函子与共尾子范畴}\etiquette[8.1]{I.sec8.1}

\begin{enonce*}{定义 8.1.1}\etiquette[8.1.1]{I.def8.1.1}
设
\hypertarget{sga4.exp.I.equation.8.1.1.1}{}$$
\displaylines{\rlap{$(8.1.1.1)$}\hfill \varphi:I\longrightarrow
I'\hfill}
$$
为一个函子. 如果对每个范畴 $C$ 和每个函子 $u:I'\rightarrow C$, 将
$\varinjlim u$ 与 $\varinjlim u\circ\varphi$ 视为
$\Hhom(G(𝒱\text{-}\Ens))^\circ$ 的对象
(\hyperref[I.def2.1]{2.1}, \hyperref[I.sec2.3.1]{2.3.1})
(这里 $𝒱$ 是使 $I$, $I'\in 𝒱$ 且 $C$ 为 $𝒱$-范畴的宇宙), 则典范态射
\hypertarget{sga4.exp.I.equation.8.1.1.2}{}$$
\displaylines{\rlap{$(8.1.1.2)$}\hfill \varinjlim
  u\circ\varphi\longrightarrow \varinjlim u\hfill}
$$
是同构, 就称函子 $\varphi$ 是共尾的 (cofinal). 当 $\varphi$ 是
$I'$ 的一个子范畴的包含函子时, 如果 $\varphi$ 共尾, 就称子范畴 $I$ 共尾.
\end{enonce*}

\setcounter{subsubsection}{1}
\subsubsection{}\etiquette[8.1.2]{I.sec8.1.2}
注意, 对给定的 $\varphi$ 和 $u$, \hyperlink{sga4.exp.I.equation.8.1.1.2}{(8.1.1.2)}
是否为双射并不依赖宇宙 $𝒱$ 的选择.

回顾 \hyperlink{sga4.exp.I.equation.8.1.1.2}{(8.1.1.2)} 中各项的定义
(\hyperref[I.def2.1]{2.1}, \hyperref[I.sec2.3.1]{2.3.1}), 可见称
$\varphi$ 共尾, 也就是说: 对每个使 $I$ 与 $I'$ 为 $𝒱$-小的宇宙 $𝒱$
以及每个函子
$$
v:{I'}^\circ\longrightarrow (𝒱\text{-}\Ens){\quoi},
$$
典范同态
\hypertarget{sga4.exp.I.equation.8.1.2.1}{}$$
\displaylines{\rlap{$(8.1.2.1)$}\hfill \varprojlim v\longrightarrow
  \varprojlim v\circ\varphi\hfill}
$$
是同构, 即双射. 为说明此条件确实必要, 只须令 $v=v^\circ$, 并注意到它也表示:
当 $C=(𝒱\text{-}\Ens)^\circ$ 时,
\hyperref[I.def8.1.1]{定义 8.1.1} 的条件成立
(这\pageoriginale{}意味着 \hyperlink{sga4.exp.I.equation.8.1.1.2}{(8.1.1.2)}
中所考虑的归纳极限在 $C$ 中可表).

\setcounter{paragraph}{1}
\paragraph{}\etiquette[8.1.2.2]{I.sec8.1.2.2}
由定义立即得出: \emph{两个共尾函子的复合仍共尾.}

\begin{enonce*}{命题 8.1.3}\etiquette[8.1.3]{I.prop8.1.3}
设 $\varphi:I\rightarrow I'$ 为函子.
\begin{enumerate}
\renewcommand{\theenumi}{\alph{enumi}}
\renewcommand{\labelenumi}{\theenumi)}
\item $\varphi$ 共尾的一个必要条件是 $\varphi$ 满足下述条件:

\begin{itemize}
\item[F 1)] 对 $I'$ 的每个对象 $i'$, 存在 $I$ 的对象 $i$, 使
$\varphi(i)$ 是 $i'$ 的一个上界 (即 $\Hom(i',\varphi(i))\neq \phi$).
\end{itemize}

\item 设 $I$ 滤过. $\varphi$ 共尾的充要条件是 $\varphi$ 满足条件
{\rm (F 1)} 及下述条件:
\begin{itemize}
\item[F 2)] 对 $I$ 的每个对象 $i$ 以及 $I'$ 中每个以 $\varphi(i)$ 为终点的双箭头
$\xymatrix{ i'\ar@<2pt>[r]^{f',g'}\ar@<-2pt>[r]& \varphi(i) }$,
存在 $I$ 中的箭头 $h:i\rightarrow j$, 使 $\varphi(h)f'=\varphi(h)g'$.
\end{itemize}
此外, 如果 $\varphi$ 共尾, 则 $I'$ 滤过.

\item 设 $I'$ 滤过且 $\varphi$ 全忠实. $\varphi$ 共尾的充要条件是它满足
{\rm a)} 中的条件 {\rm F 1)}. 此时 $I$ 滤过.
\end{enumerate}
\end{enonce*}

\noindent \emph{证明.} 对 a) 与 b) 中的必要性, 只须对 $u$ 为典范包含函子
(\hyperref[I.sec1.3.3]{1.3.3}) 的情形应用定义 \hyperref[I.def8.1.1]{8.1.1}:
$$
u:I'\hookrightarrow \widehat{I}'
$$
(选取一个使 $I$ 与 $I'$ 为 $𝒰$-小的宇宙 $𝒰$).
于是可将 \hyperlink{sga4.exp.I.equation.8.1.1.2}{(8.1.1.2)} 解释为
$\widehat{I}'$ 中的一条箭头 (\hyperref[I.prop3.1]{3.1}), 其终点是
$I'$ 上的\emph{终}\pageoriginale{}函子 (\hyperref[I.prop3.4]{3.4}).
条件 F 1) 显然表示这条箭头是 $\widehat{I}'$ 中的满态射, 即在每个变元处均满射
(注意 $\widehat{I}'$ 中的归纳极限是“逐变元”计算的). 这证明了 a).
现设 $I$ 滤过且 $\varphi$ 共尾. 则 $I'$ 滤过: 事实上,
$I'$ 的任意两个对象有第三个对象作为共同上界这一条件, 立即由 $I$ 的相同条件及 F 1) 得出;
只须再证 \hyperref[I.def2.7]{2.7} 中的条件 PS 2), 即对 $I'$ 中每个双箭头
$$
f',g':i'\rightrightarrows j'
$$
存在 $I'$ 中的箭头 $h':j'\rightarrow k'$, 使 $h'f'=h'g'$.
由 F~1), 可设 $k'$ 形如 $\varphi(i)$. 又因 $\varphi$ 共尾, 对 $I'$ 的每个对象 $i'$,
集合 $\varinjlim_i\Hom(i',\varphi(i))$ 仅有一个元素.
因此, (Ens) 中\emph{滤过}归纳极限的通常描述
(\hyperref[I.lem2.8.1]{2.8.1}) 表明, 存在 $I$ 中的箭头 $h:i\rightarrow j$,
使 $\varphi(h)f'=\varphi(h)g'$. 这就证明了 $I'$ 滤过, 同时也证明了条件 F~2).

为证明 b) 中的条件足以保证 $\varphi$ 共尾, 使用定义的形式
\hyperref[I.sec8.1.2]{8.1.2}. 立即可见, 条件 F 1) 蕴含
\hyperlink{sga4.exp.I.equation.8.1.2.1}{(8.1.2.1)} 为单态射, 即在每个变元处均单射;
而条件 F 2) (连同 F 1)) 保证它是双射
(注意 $\widehat{{I'}^\circ}$ 中的投射极限逐变元计算). 这证明了 b).

最后,\pageoriginale{}若 $I'$ 滤过且 $\varphi$ 全忠实, 则条件 F 1) 显然蕴含 $I$ 滤过,
并蕴含条件 F 2). 因而 c) 由 a) 与 b) 得出.

\setcounter{subsubsection}{3}
\subsubsection{}\etiquette[8.1.4]{I.sec8.1.4}
当 $I'$ 为滤过范畴而 $I$ 为 $I'$ 的满子范畴时, 由
\hyperref[I.prop8.1.3]{8.1.3} c) 可见, $I$ 在 $I'$ 中共尾这一条件仅依赖于
预序集 $OB\,I'$ 的子集 $\Ob I$ (其预序关系为
“$x\leq y$” $\Leftrightarrow$ “$\Hom(x,y)\neq ∅$”).
如果 $J'$ 为预序集而 $J$ 为 $J'$ 的子集, 则当 $J'$ 的每个元素都以 $J$ 的某个元素为上界时,
我们有时称 $J$ 是 $J'$ 的\emph{共尾子集}.
因而当 $J'$ 滤过时, 这就表示相应范畴的包含函子 $𝒥\hookrightarrow 𝒥'$ 共尾.

\subsubsection{}\etiquette[8.1.5]{I.sec8.1.5}
下文仅在范畴 $I$ 与 $I'$ 滤过时使用共尾函子的概念.
经典处理中甚至只考虑预序集对应的范畴
(即给定起点与终点时至多有一条箭头的范畴).
但这种限制在应用中并不方便, 因为许多应用中自然出现的滤过范畴\emph{并非}有序范畴.
不过, 下述归于 P. \textsc{Deligne} 的结果表明,
这两种观点之间并无本质区别:

\begin{enonce*}{命题 8.1.6}\etiquette[8.1.6]{I.prop8.1.6}
设\pageoriginale $I$ 为小滤过范畴. 则存在小有序集 $E$ 以及共尾函子
$\varphi:ℰ\rightarrow I$, 其中 $ℰ$ 表示 $E$ 对应的范畴.
\end{enonce*}

先设预序集 $\Ob I$ 没有最大元.
$I$ 的子图是由 $F\ell(I)$ 的子集 $F$ 和 $\Ob I$ 的子集 $\mathcal{O}$ 组成的二元组
$D=(\mathcal{O},F)$, 且对每个 $f\in F$, $f$ 的起点和终点均属于 $\mathcal{O}$.
如果对每个 $x\in D'$, $D$ 中起点为 $x$、终点为 $y$ 的箭头所成集合
$\Hom(x,e)\cap F$ 恰有一个元素 $f_x$, 对 $D$ 的每条箭头
$g:x\rightarrow y$ 均有 $f_x=f_y\circ g$, 并且 $f_e=\id_e$,
则称 $\mathcal{O}$ 的元素 $e$ 为子图 $D$ 的一个\emph{终对象}.
令 $E$ 为 $I$ 的所有具有\emph{唯一}终元素 $\varphi(D)$ 的\emph{有限}子图 $D$ 所成的集合,
并按包含关系将 $E$ 排序. 若 $D$, $D'$ 为 $E$ 的两个元素且 $D\subset D'$,
则 $D'$ 中存在唯一一条起点为 $\varphi(D)$、终点为 $\varphi(D')$ 的箭头
$\varphi(D',D)$. 若 $D\subset D'\subset D''$ 为 $E$ 中的包含关系,
显然有 $\varphi(D'',D)=\varphi(D'',D')\varphi(D',D)$; 又有
$\varphi(D,D)=\id_{\varphi(D)}$. 由此得到函子 $\varphi:ℰ\rightarrow I$,
余下须证 $E$ 滤过且 $\varphi$ 共尾. 由 \hyperref[I.prop8.1.3]{8.1.3} b),
须作三项验证:
\begin{enumerate}
\renewcommand{\labelenumi}{\theenumi)}
\item \hyperref[I.prop8.1.3]{8.1.3} 的条件 F 1): 对每个 $i\in I$,
取 $D$ 为仅由对象 $i$ 及其恒等箭头组成的 $I$ 的子图; 则 $i\leq \varphi(D)=i$.
此条件 F 1) 也表明 $E\neq ∅$.

\item 条件 F 2): 对每个 $i\in \Ob I$、每个 $D\in E$ 以及每个双箭头
$i\rightrightarrows \varphi(D)$, 要找到包含 $D$ 的图 $D'\in E$,
使 $\varphi(D',D):\varphi(D)\rightarrow \varphi(D')$ \emph{等化}该双箭头.
由于 $I$ 滤过,\pageoriginale{}存在 $I$ 的对象 $j$ 以及箭头
$f:\varphi(D)\rightarrow j$ 等化该双箭头. 由于 $I$ 没有最大元,
可设 $j$ 为 $\varphi(D)$ 的\emph{严格}上界, 因而它不同于 $D$ 的任何其他对象.
令 $D'$ 为 $I$ 的如下子图: 其对象集为 $\mathcal{O}\cup\{j\}$,
其箭头集由 $D$ 的箭头、各复合
$x\xrightarrow{f_x}\varphi(D)\xrightarrow{f}j$
以及 $\id_j$ 组成, 这里 $x$ 为 $D$ 的对象, $f_x$ 为 $D$ 中起点为 $x$、终点为
$\varphi(D)$ 的唯一箭头. 显然 $D'$ 是 $I$ 的有限子图, 且以 $j$ 为唯一终对象,
因而 $D'\in E$, 而且 $D'$ 满足所需条件.

\renewcommand{\labelenumi}{$3^\circ$}
\item 两个子图 $D$, $D'\subset E$ 包含于同一个 $D''\in E$.
事实上, 可取 $\varphi(D)$ 与 $\varphi(D')$ 的一个严格上界 $j$,
\[
\xymatrix@C=1.5cm@R=1cm{
\varphi(D)\ar[rd]^f && \\
           &j &\\
\varphi(D')\ar[ru]_{f'}&\quad , }
\]
并令 $D''$ 为如下子图: 其对象集是 $D$ 与 $D'$ 的对象集同 $\{j\}$ 的并,
其箭头集是 $D$ 与 $D'$ 的箭头集、复合 $f\circ f_x$ ($x$ 为 $D$ 的对象)
与 $f'\circ f'_{x'}$ ($x'$ 为 $D'$ 的对象) 所成集合以及 $\{\id_j\}$ 的并.
这样确实得到 $E$ 的一个有限子图. 下面说明, 对适当的 $g:j\rightarrow j'$,
将 $j$, $f$, $f'$ 换为 $j'$, $gf$, $gf'$ 后, 可使 $j$ 成为 $D'$ 的终对象.
等价地, 对同时为 $D$ 和 $D'$ 对象的每个 $x$, 必须有 $ff_x=f'f'_{x'}$.
这归结为要用某个 $g:j\rightarrow j'$ 等化\emph{有限}个双箭头,\pageoriginale
而 $I$ 滤过保证这是可能的.
\end{enumerate}

这就完成了所考虑情形的证明. 引入自然数有序集 $\NN$ 对应的滤过范畴 $\NN$,
并注意范畴 $\NN\times I$ 滤过、没有最大元, 而且投影
$\NN\times I\rightarrow I$ 为共尾函子, 就可将一般情形归结为上述情形.

\begin{enonce*}{推论 8.1.7}\etiquette[8.1.7]{I.cor8.1.7}
设 $I$ 为 $𝒰$-范畴. 存在小滤过有序集 $E$ 以及共尾函子 $ℰ\rightarrow I$
的充要条件是 $I$ 滤过且 $\Ob I$ 有一个小共尾子集 $I'$
(\hyperref[I.sec8.1.4]{8.1.4}).
\end{enonce*}

必要性由 \hyperref[I.prop8.1.3]{8.1.3} a) 得出;
充分性由对 $I'$ 定义的 $I$ 的满子范畴应用 \hyperref[I.prop8.1.6]{8.1.6} 得出.

\begin{enonce*}{定义 8.1.8}\etiquette[8.1.8]{I.def8.1.8}
设 $I$ 为滤过范畴. 如果 $I$ 为 $𝒰$-范畴且满足
\hyperref[I.cor8.1.7]{8.1.7} 的各等价条件, 就称 $I$ 本质小.
\end{enonce*}

\setcounter{subsection}{1}

% END INCLUDED SGA4_English_current_buildable_source_20260906/source/01/drafts/SGA4_I_LINES_2669_2936_English.texfrag


% BEGIN INCLUDED SGA4_English_current_buildable_source_20260906/source/01/drafts/SGA4_I_LINES_2937_3118_English.texfrag
\subsection{ind-对象与 ind-可表函子}\etiquette[8.2]{I.sec8.2}

\subsubsection{}\etiquette[8.2.1]{I.sec8.2.1}
下文固定一个 $𝒰$-范畴 $C$, 并始终借助典范函子
\hyperref[I.sec1.3.1]{(1.3.1)}
\hypertarget{sga4.exp.I.equation.8.2.1.1}{}$$
\displaylines{\rlap{$(8.2.1.1)$}\hfill h:C\hookrightarrow
  \widehat{C}{\quoi}.\hfill}
$$
将其视为嵌入 $𝒰$-预层范畴 $\widehat{C}$ 中.
$C$ 的一个 \emph{ind-对象} (或 $C$ 的一个\emph{归纳系}) 是任意一个函子
\hypertarget{sga4.exp.I.equation.8.2.1.2}{}$$
\displaylines{\rlap{$(8.2.1.2)$}\hfill \varphi:I\longrightarrow
  C{\quoi},\hfill}
$$
其中\pageoriginale $I$ 为\emph{滤过}范畴 (\hyperref[I.def2.7]{2.7}),
称为该 ind-对象的\emph{指标范畴}. (除此以外, $I$ 任意; 尤其不要求 $I$ 为 $𝒰$-范畴.)
通常用带指标的记号 $(X_i)_{i\in \ob I}$, 甚至 (再作一次记号滥用)
$(X_i)_{i\in I}$、$(X_i)_i$ 或 $(X_i)$ 表示 ind-对象 $\varphi$ 会比较方便,
其中对 $i\in \ob I$, 有 $X_i=\varphi(i)$.
这种记号与经典中用于以滤过有序集为指标的归纳系的记号
(同样不提及过渡态射) 所作的滥用程度完全相同.
各 $X_i$ 称为 ind-对象 $𝒳$ 的\emph{分量对象}.
注意, 在这里考虑的一般情形中, 归纳系的“\emph{过渡态射}” $\varphi(f)$
以 $I$ 的箭头集 $F\ell I$ 为指标, 而后者一般不能等同于
$\Ob I\times\Ob I$ 的子集; 换言之, 对给定的 $i$ 和 $j$, 若 $j$ 为 $i$ 的上界,
从 $X_i$ 到 $X_j$ 可能存在多个过渡态射.

\subsubsection{}\etiquette[8.2.2]{I.sec8.2.2}
$C$ 的最有用的 ind-对象 $𝒳=(X_i)_{i\in I}$ 是指标范畴 $I$ 本质小的那些
\hyperref[I.def8.1.8]{(8.1.8)}; 这样的 ind-对象称为\emph{本质小}的.
若 $(X_i)_{i\in I}$ 为这样的对象, 则由小归纳极限在 $\widehat{C}$ 中可表
\hyperref[I.prop3.1]{(3.1)}, 可考虑
\hypertarget{sga4.exp.I.equation.8.2.2.1}{}$$
\displaylines{\rlap{$(8.2.2.1)$}\hfill
L(𝒳)\mathrel{\mathop=^{\rm def}}
  \varinjlim h\cdot 𝒳 = \varinjlim_i X_i \in \ob
  \widehat{C}{\quoi},\hfill}
$$
其中最后一个极限是\emph{在} $\widehat{C}$ \emph{中}取得的.
因而 $L(𝒳)$ 是 $C$ \pageoriginale{}上的预层
\hypertarget{sga4.exp.I.equation.8.2.2.2}{}$$
\displaylines{\rlap{$(8.2.2.2)$}\hfill L(𝒳):Y\longmapsto
  \varinjlim_i \Hom(Y,X_i)\hfill}
$$
此预层称为\emph{由 $C$ 的 ind-对象 $𝒳=(X_i)_{i\in I}$ ind-表示的预层}.
$C$ 上的预层 $F$ 如同构于某个由本质小 ind-对象 ind-表示的预层 $L(𝒳)$,
就称为 \emph{ind-可表预层}. 此外, 由 \hyperref[I.prop8.1.6]{8.1.6} 可知,
在此定义中可设 $I$ 为小范畴, 甚至可设 $I$ 对应于某个有序集 $\in 𝒰$.

\subsubsection{}\etiquette[8.2.3]{I.sec8.2.3}
考虑由函子 $\varphi:I\rightarrow C$ 给出的 ind-对象 $𝒳=(X_i)_{i\in I}$,
并设
$$
u:I'\longrightarrow I
$$
为函子, 其中 $I'$ 为另一个滤过范畴. 则函子 $\varphi\circ u$ 是
$C$ 的一个以 $I'$ 为指标范畴的 ind-对象 $𝒳'$; 其带指标的记号为
$(X_{u(i')})_{i'\in I'}$. 称它为借助函子 $u$ \emph{变换指标范畴,
从 $(X_i)_{i\in I}$ 得到的} $C$ 的 ind-对象.
如果 $I$ 与 $I'$ (即 $𝒳$ 与 $𝒳'$) 本质小, 就得到分别由 $𝒳'$ 与 $𝒳$
ind-表示的预层之间的典范态射
\hypertarget{sga4.exp.I.equation.8.2.3.1}{}$$
\displaylines{\rlap{$(8.2.3.1)$}\hfill
  L(𝒳')=\varinjlim_{i'}X_{u(i')}\longrightarrow L(𝒳)=\varinjlim_i
  X_i\hfill}
$$
当 $u$ 为共尾函子 (\hyperref[I.def8.1.1]{8.1.1}) 时,
$𝒳$ 本质小当且仅当 $𝒳'$ 本质小 \hyperref[I.cor8.1.7]{(8.1.7)},
且上述同态 \hyperlink{sga4.exp.I.equation.8.2.3.1}{(8.2.3.1)} 是同构
\hypertarget{sga4.exp.I.equation.8.2.3.2}{}$$
\displaylines{\rlap{$(8.2.3.2)$}\hfill L(𝒳')\simeq
  L(𝒳){\quoi}.\hfill}
$$

\subsubsection{}\etiquette[8.2.4]{I.sec8.2.4}
设\pageoriginale
\hypertarget{sga4.exp.I.equation.8.2.4.1}{}$$
\displaylines{\rlap{$(8.2.4.1)$}\hfill 𝒳=(X_i)_{i\in
    I},𝒴=(Y_j)_{j\in I}\hfill}
$$
为 $C$ 的两个本质小 ind-对象, 分别以 (滤过且本质小的) 范畴 $I$ 与 $J$ 为指标.
从 ind-对象 $𝒳$ 到 ind-对象 $𝒴$ 的一个\emph{态射}, 就是它们所 ind-表示的预层之间的任意态射
\hypertarget{sga4.exp.I.equation.8.2.4.2}{}$$
\displaylines{\rlap{$(8.2.4.2)$}\hfill L(𝒳)\longrightarrow
  L(𝒴)\hfill}
$$
令
\hypertarget{sga4.exp.I.equation.8.2.4.3}{}$$
\displaylines{\rlap{$(8.2.4.3)$}\hfill
\Hom_{\text{ind}-\ob}(𝒳,𝒴)=\Hom_{\widehat{C}}(L(𝒳),L(𝒴)){\quoi}.
   \hfill}
$$
(在不会混淆时, 省略 $\Hom$ 的下标 ind-ob.)
$𝒞$ 的 ind-对象态射的复合, 定义为它们所 ind-表示的预层的态射的复合.
若 $𝒱\supset 𝒰$ 为包含 $𝒰$ 的宇宙, 并且只考虑属于 $𝒱$ 的本质小 ind-对象,
则这些对象构成一个范畴的对象集; 其箭头集由三元组 $(𝒳,𝒴,u)$ 组成,
其中 $𝒳$ 和 $𝒴$ 为本质小 ind-对象 $\in 𝒱$,
$u:𝒳\rightarrow 𝒴$ 为 ind-对象的态射, 即态射 $L(𝒳)\rightarrow L(𝒴)$.
这样得到的范畴记为
\hypertarget{sga4.exp.I.equation.8.2.4.4}{}$$
\displaylines{\rlap{$(8.2.4.4)$}\hfill
  \Ind_{𝒱}(C,𝒰){\quoi}.\hfill}
$$
当 $𝒱=𝒰$ 时, 简记为
\hypertarget{sga4.exp.I.equation.8.2.4.5}{}$$
\displaylines{\rlap{$(8.2.4.5)$}\hfill \Ind_{𝒰}(C)\text{ 或
}
  \Ind(C){\quoi},\hfill}
$$
并称为相对于 $𝒰$ 的 \emph{$C$ 的 ind-对象范畴}; 在不会混淆时省略 $𝒰$.

显然,\pageoriginale{}若 $𝒱\subset 𝒱'$ 为两个包含 $𝒰$ 的宇宙,
则 $\Ind_{𝒱}(C,𝒰)$ 为 $\Ind_{𝒱'}(,𝒰)$ 的满子范畴;
由 \hyperref[I.sec8.2.3]{8.2.3} 得出, 包含函子是\emph{范畴等价}:
\hypertarget{sga4.exp.I.equation.8.2.4.6}{}$$
\displaylines{\rlap{$(8.2.4.6)$}\hfill
\Ind_{𝒱}(C,𝒰)\xrightarrow{\approx}\Ind_{𝒱'}(C,𝒰){\quoi}.\hfill}
$$
这尤其表明, \emph{$C$ 的每个本质小 ind-对象都确定 $\Ind(C)$ 的一个元素,
精确到唯一同构.} 这说明了把 $C$ 的每个本质小 ind-对象等同于
$\Ind(C)$ 的一个对象这一颇为常见的用语滥用的合理性.

由定义立即可知, 对每个宇宙 $𝒱$, 存在典范函子
\hypertarget{sga4.exp.I.equation.8.2.4.7}{}$$
\displaylines{\rlap{$(8.2.4.7)$}\hfill
  L:\Ind_{𝒱}(C,𝒰)\longrightarrow \widehat{C}\hfill}
$$
且它全忠实. 由 \hyperlink{sga4.exp.I.equation.8.2.4.6}{(8.2.4.6)},
显然只要知道其中一个函子, 尤其是知道典范函子
\hypertarget{sga4.exp.I.equation.8.2.4.8}{}$$
\displaylines{\rlap{$(8.2.4.8)$}\hfill L:\Ind(C)\longrightarrow
  \widehat{C}{\quoi}.\hfill}
$$
这些函子便被确定, 精确到唯一同构.
由于此函子全忠实, 常用它将左端的一个对象与其所 ind-表示的函子等同,
甚至将 $\Ind(C)$ 与它在 $\widehat{C}$ 中的本质像
(由 ind-可表预层组成) 等同. 但须注意, 相比借助函子
$h$ \hyperlink{sga4.exp.I.equation.8.2.1.1}{(8.2.1.1)} 将 $C$ 等同于
$\widehat{C}$ 的满子范畴, 这种等同明显有更多不便,
因为与后者不同, 函子 \hyperlink{sga4.exp.I.equation.8.2.4.8}{(8.2.4.8)} 一般在对象上并不单射.


% END INCLUDED SGA4_English_current_buildable_source_20260906/source/01/drafts/SGA4_I_LINES_2937_3118_English.texfrag


% BEGIN INCLUDED SGA4_English_current_buildable_source_20260906/source/01/drafts/SGA4_I_LINES_3119_3227_English.texfrag
\subsubsection{}\etiquette[8.2.5]{I.sec8.2.5}
下面用所考虑的 ind-对象的带指标表达式
\hyperlink{sga4.exp.I.equation.8.2.4.1}{(8.2.4.1)} 将定义
\hyperlink{sga4.exp.I.equation.8.2.4.3}{(8.2.4.3)} 明确写出. 考虑定义
\hyperlink{sga4.exp.I.equation.8.2.2.1}{(8.2.2.1)}, 得到\pageoriginale
$$
\Hom(𝒳,𝒴)\simeq \varprojlim_i
\Hom(X_i,𝒴)\mathrel{\mathop=^{\rm dfn}} \varprojlim_i
\Hom(X_i, L(𝒴))=\varprojlim_i\varinjlim_j \Hom_C(X_i,
Y_j){\quoi},
$$
其中最后一个等式由对 $𝒴$ 应用
\hyperlink{sga4.exp.I.equation.8.2.2.2}{(8.2.2.2)} 得出.
于是得到典范双射
\hypertarget{sga4.exp.I.equation.8.2.5.1}{}$$
\displaylines{\rlap{$(8.2.5.1)$}\hfill \Hom((X_i)_{i\in
I},(Y_j)_{j\in
    J})\simeq \varprojlim_{i\in I} \varinjlim_{j\in
    J}\Hom_C(X_i,Y_j){\quoi}.\hfill}
$$
将 ind-对象态射的复合用此公式明确表达出来, 留待读者完成.
注意, 由此公式立即得出, ind-对象的同态集 $\Hom(𝒳,𝒴)$ 是
$𝒰$-\emph{小}的, 即各范畴 \hyperlink{sga4.exp.I.equation.8.2.4.4}{(8.2.4.4)}
都是 $𝒰$-范畴. 等价地 (由 \hyperlink{sga4.exp.I.equation.8.2.4.6}{(8.2.4.6)}),
\emph{范畴 $\Ind(C)$ 是 $𝒰$-范畴}, 即当 $I$, $J\in 𝒰$ 时,
\hyperlink{sga4.exp.I.equation.8.2.5.1}{(8.2.5.1)} 的右端为小集.
这立即来自 $C$ 为 $𝒰$-范畴这一事实, 因而各集合 $\Hom_C(X_i,Y_j)$ 都小.

\subsubsection{}\etiquette[8.2.6]{I.sec8.2.6}
设 $I$ 为本质小滤过范畴. 由
\hyperlink{sga4.exp.I.equation.8.2.5.1}{(8.2.5.1)}, 或由函子
$I\rightarrow \widehat{C}$ 的归纳极限对该函子的函子性, 得出有典范函子
\hypertarget{sga4.exp.I.equation.8.2.6.1}{}$$
\displaylines{\rlap{$(8.2.6.1)$}\hfill \Hhom(I,C)\longrightarrow
  \Ind(C){\quoi}.\hfill}
$$
注意, 此函子一般\emph{并不}全忠实, 甚至并不忠实.

\begin{enonce*}[remark]{注 8.2.7}\etiquette[8.2.7]{I.sec8.2.7}
\hyperref[I.sec8.2.4]{8.2.4} 中的\pageoriginale{}定义明确了两个 (本质小)
ind-对象的\emph{同构}概念, 这也就是它们所 ind-表示的预层的同构.
注意, 相比形如 $\Hhom(I,C)$ 的范畴中的同构概念, 这是一个很弱的同构概念.
因此, 可对 ind-对象提出的许多颇为自然的条件
(例如要求分量属于给定的严格满子范畴, 或要求它严格 (下文 \hyperref[I.sec8.12]{8.12}))
\emph{并不}在同构下保持不变.
因此, 若希望使用在 ind-对象同构下不变的概念, 就须将所讨论的概念“饱和化”:
从由满足该条件的对象所成的 $\Ind(C)$ 的子范畴,
过渡到它所生成的 $\Ind(C)$ 的严格满子范畴.
例如, 参看下文引入的本质常值 ind-对象概念 (\hyperref[I.sec8.4]{8.4}).
\end{enonce*}

\begin{enonce*}[remark]{习题 8.2.8}\etiquette[8.2.8]{I.exer8.2.8}
设 $𝒰\subset 𝒱$ 为两个宇宙, $C$ 为属于 $𝒱$ 的 $𝒰$-范畴.
用 $\SysInd(C)$ 表示下述范畴:
\begin{enumerate}
\renewcommand{\theenumi}{\alph{enumi}}
\renewcommand{\labelenumi}{\theenumi)}
\item $\SysInd(C)$ 的对象是函子 $\varphi:I\rightarrow C$,
其中 $I$ 为属于 $𝒱$ 的本质 $𝒰$-小滤过范畴.

\item 设 $(I,\varphi)$ 与 $(J,\psi)$ 为 $\SysInd(C)$ 的两个对象.
$\SysInd(C)$ 中从 $(I,\varphi)$ 到 $(J,\psi)$ 的一个态射是二元组 $(m,u)$,
其中 $m:I\rightarrow J$ 为函子, $u:\varphi\rightarrow \psi\circ m$ 为函子态射.
$\SysInd(C)$ 中态射的复合按显然的方式定义.
\end{enumerate}
\end{enonce*}

设\pageoriginale $S$ 为 $\SysInd(C)$ 中满足 $m$ 共尾且 $u$ 为同构的态射
$(m,u):(I,\varphi)\rightarrow (J,\Psi)$ 所成的集合.
设 $(I,\varphi)$ 为 $\SysInd(C)$ 的对象.
$I$ 的箭头范畴 $F\ell(I)$ 通过起点和终点这两个自然函子映至 $I$.
这两个函子又由典范函子态射 $v:$ 起点 $\rightarrow$ 终点相联系.
因此 $\SysInd(C)$ 中有两个从 $(F\ell(I),\varphi\circ\text{ 起点})$
到 $(I,\varphi)$ 的态射:
$p_1(I,\varphi)=(\text{起点, } \id)$ 及
$p_2(I,\varphi)=(\mbox{终点, } \varphi^\ast v:\varphi\circ
\text{起点 }\rightarrow \varphi\circ{\text{ 终点}})$.
令 $p:\SysInd(C)\rightarrow\Ind(C)$ 为显然的函子. 设 $B$ 为范畴.
证明函子 $F\mapsto F\circ p$
$$
\Hhom(\Ind(C),B)\longrightarrow \Hhom(\SysInd(C),B)
$$
全忠实, 且函子 $G:\SysInd(C)\rightarrow B$ 属于它的本质像,
当且仅当它具有以下两个性质:
\begin{enumerate}
\renewcommand{\labelenumi}{\theenumi)}
\item 对每个 $s\in S$, $F(s)$ 是 $B$ 的同构.

\item 对 $\SysInd(C)$ 的每个对象 $(I,\varphi)$,
$F(p_1(I,\varphi))=F(p_2(I,\varphi))$.
\end{enumerate}


% END INCLUDED SGA4_English_current_buildable_source_20260906/source/01/drafts/SGA4_I_LINES_3119_3227_English.texfrag


% BEGIN INCLUDED SGA4_English_current_buildable_source_20260906/source/01/drafts/SGA4_I_LINES_3228_3408_English.texfrag
\subsection{ind-可表函子的刻画}\etiquette[8.3]{I.sec8.3}

\begin{enonce*}{命题 8.3.1}\etiquette[8.3.1]{I.prop8.3.1}
设 $F$ 为 $C$ 上的 ind-可表预层. 则 $F$ 左正合, 即
(\hyperref[I.cor2.3.2]{2.3.2}) 对每个有限范畴 $J$ 和每个使
$\varinjlim \varphi$ 可表 (即 $\varprojlim \varphi^\circ$ 可表) 的函子
$\varphi:J\rightarrow C$, 典范映射
$$
F(\varinjlim \varphi)\longrightarrow \varprojlim F\circ\varphi
$$
为双射.
\end{enonce*}

事实上,\pageoriginale{}可表预层显然左正合, 由 \hyperref[I.prop2.8]{2.8} 立即得出,
这类函子的任意滤过归纳极限也左正合, \cqfd

\setcounter{subsubsection}{1}
\subsubsection{}\etiquette[8.3.2]{I.sec8.3.2}
给定 $C$ 上的预层 $F$, 我们将常用到范畴
\hypertarget{sga4.exp.I.equation.8.3.2.1}{}$$
\displaylines{\rlap{$(8.3.2.1)$}\hfill C_{/F}\hookrightarrow
\widehat{C}_{/F}{\quoi},\hfill}
$$
它表示后一个范畴中由\emph{起点属于} $C$ 的箭头 $X\rightarrow F$ 组成的满子范畴.
由 \hyperref[I.prop1.4]{1.4}, 此范畴的对象可等同于二元组 $(X,u)$,
其中 $X$ 为 $C$ 的对象, $u\in F(X)$.
从 $(X,u)$ 到 $(X',u')$ 的态射遂解释为 $C$ 中满足 $F(f)(u')=u$ 的箭头
$f:X\rightarrow X'$. 从 $\widehat{C}_{/F}$ 到 $\widehat{C}$ 的
“遗忘终点”函子诱导函子 (称为\emph{典范函子})
\hypertarget{sga4.exp.I.equation.8.3.2.2}{}$$
\displaylines{\rlap{$(8.3.2.2)$}\hfill C_{/F}\longrightarrow
C{\quoi},\hfill}
$$
它已在 \hyperref[I.prop3.4]{3.4} 中考虑过; 那里证明了此函子的归纳极限在
$\widehat{C}$ 中存在 (对 $C_{/F}$ 不加任何小性条件!), 且典范同构于 $F$.


\setcounter{paragraph}{2}
\paragraph{}\etiquette[8.3.2.3]{I.sec8.3.2.3}
注意, 若有限归纳极限在 $C$ 中可表, 且 $F$ 左正合
(即将这些极限映为 $(\Ens)$ 中的有限投射极限), 则 $C_{/F}$ 中也如此,
因而更有 (\hyperref[I.exem2.7.1]{2.7.1}) $C_{/F}$ 滤过.
更一般地, 若两个对象的和以及平行箭头对的余核在 $C$ 中可表,
且 $F$ 将它们分别映为积和核, 则\pageoriginale $C_{/F}$
对同类有限归纳极限封闭, 因而它滤过当且仅当它非空,
即当且仅当函子 $F$ 不是取值为 $∅$ 的常值函子.

\begin{enonce*}{定理 8.3.3}\etiquette[8.3.3]{I.them8.3.3}
设 $C$ 为 $𝒰$-范畴, $F\in \ob \widehat{C}$,
$F:C^\circ\rightarrow (𝒰\text{-}\Ens)$ 为 $C$ 上的 $𝒰$-预层.
下述条件等价:
\begin{itemize}
\item[(i)] $F$ ind-可表 (\hyperref[I.sec8.2.2]{8.2.2}).

\item[(ii)] 范畴 $C_{/F}$ (\hyperref[I.sec8.3.2]{8.3.2}) 滤过且本质小.

\item[(iii)] (若有限归纳极限在 $C$ 中可表.) 函子 $F$ 左正合,
且 $\ob C_{/F}$ 有一个小共尾子集 (\hyperref[I.sec8.1.4]{8.1.4}).

\item[(iii bis)] (若两个对象的和及平行箭头对的余核在 $C$ 中可表.)
函子 $F$ 将 $C$ 中两个对象的和映为积, 将 $C$ 中平行箭头对的余核映为核;
范畴 $C_{/F}$ 非空, 即函子 $F$ 不是取值为 $∅$ 的常值函子;
最后, 存在 $C_{/F}$ 的一个小对象族, 使 $C_{/F}$ 的每个对象 $X$
都由某个对象 $X_i$ 支配, 即有一个 $F$-态射 $X_i\rightarrow X$.

\item[(iv)] (若范畴 $C$ 等价于小范畴.) 范畴 $C_{/F}$ 滤过.

\item[(v)] (若范畴 $C$ 等价于小范畴, 且滤过归纳极限在其中可表.)
函子 $F$ 左正合.
\end{itemize}
\end{enonce*}

\noindent \emph{证明.} (i)\pageoriginale $\Rightarrow$
(ii). 设 $F$ 由 $(X_i)_{i\in I}$ ind-表示, 其中 $I$ 小.
先证 $C_{/F}$ 滤过. 设 $X$, $X'$ 为 $C_{/F}$ 的两个对象,
即配备了态射 $X\rightarrow F$, $X'\rightarrow F$ 的 $C$ 的对象;
要证它们在 $C_{/F}$ 中有共同上界.
态射 $X\rightarrow F$, $X'\rightarrow F$ 来自态射
$X\rightarrow X_i$, $X'\rightarrow X'_i$;
将 $i$, $i'$ 换成它们在 $\Ob I$ 中的共同上界后, 可设 $i=i'$,
并取配备典范态射 $X_i\rightarrow F$ 的对象 $X_i$ 为 $X$, $X'$ 的共同上界.
现设 $\xymatrix{
X\ar@<2pt>[r]^{f,g}\ar@<-2pt>[r]& X' \ar@<2pt>[r]^{h} & F }$
为 $C_{/F}$ 中的一个平行箭头对; 要证它由 $C_{/F}$ 中的一条箭头
$X'\rightarrow X''$ 等化.
现有 $h:X'\rightarrow F$ 由某个态射 $h_i:X'\rightarrow X_i$ 给出,
条件 $hf=hg$ 表示存在 $I$ 中的 $\alpha:i\rightarrow j$,
使 $\varphi(\alpha)(h_if)=\varphi(\alpha)(h_ig)$,
即将 $h_i$ 换为 $\varphi(\alpha)h_i$ 后, 此态射等化 $f$ 和 $g$.
这证明 $C_{/F}$ 滤过. 随即显然它本质小,
因为对象 $(X_i\rightarrow F)_{\in \Ob I}$ 构成 $\Ob C_{/F}$ 中的小共尾族.

(ii) $\Rightarrow$ (i). 令 $I=C_{/F}$, 并考虑典范函子
\hyperlink{sga4.exp.I.equation.8.3.2.1}{(8.3.2.1)}
$$
\varphi:I=C_{/F}\longrightarrow F{\quoi}.
$$
我们知道, 此 $C$ 的 ind-对象所 ind-表示的预层是 $F$
(\hyperref[I.prop3.4]{3.4}), 故由定义 (\hyperref[I.sec8.2.2]{8.2.2}),
$F$ ind-可表.

(ii) $\Leftrightarrow$ (iii). 事实上, (ii) $\Rightarrow$ (iii),
因为 ind-可表函子左正合 \hyperref[I.prop8.3.1]{(8.3.1)};
而 (iii) $\Rightarrow$ (ii), 因为我们已注意到
(\hyperref[I.prop8.3.1]{8.3.1}), 若 $F$ 左正合, 则 $C_{/F}$ 滤过,
再\pageoriginale{}应用定义 \hyperref[I.def8.1.8]{8.1.8}, 便得此范畴本质小.

(ii) $\Leftrightarrow$ (iii bis) 的证明与 (ii) $\Leftrightarrow$ (iii) 相同.

等价关系 (ii) $\Leftrightarrow$ (iv) 及 (iii) $\Leftrightarrow$ (v) 是直接的,
因为若 $C$ 等价于小范畴, 则 $C_{/F}$ 显然也等价于小范畴,
故只要它滤过, \textit{就更}自动本质小.
这就完成了 \hyperref[I.them8.3.3]{8.3.3} 的证明.

\begin{enonce*}[remark]{注 8.3.4}\etiquette[8.3.4]{I.rem8.3.4}
设 $F:C'\rightarrow C''$ 为 $𝒰$-范畴之间的函子.
看来, 除非有限投射极限在 $C'$ 中可表, 否则左正合性
(\hyperref[I.cor2.3.2]{2.3.2}) 的概念没有多少实际意义.
当 $C''=(𝒰\text{-}\Ens)$ 且 $C'$ 为 $𝒰$-小时,
将 $C'$ 写为 $C^\circ$ (故 $C={C'}^\circ)$,
就应当认为, 在一般情形中代替左正合性的“正确概念”,
是 $F$ 的 ind-可表性 (将其视为 $C^\circ$ 上的预层;
即将 $F$ 视为 $C'$ 上的函子时的 pro-可表性,
参见 \hyperref[I.sec8.10]{8.10}).
当有限投射极限在 $C'$ 中可表, 即有限归纳极限在 $C$ 中可表时,
此概念确实与左正合性相同 (\hyperref[I.them8.3.3]{8.3.3}).
当不再假定 $C'$ 为 $𝒰$-小, 甚至不假定它等价于 $𝒰$-小范畴时,
函子 $F$ 的左正合性与 pro-可表性这两个概念便未必一致,
即使 $C'$ 中的有限投射极限可表也是如此
(参见 \hyperref[I.exer8.12.9]{8.12.9}).
事实上, 似乎在\pageoriginale{}这种情形中,
应将非 ind-可表的左正合函子视为病态对象,
而“好”对象仍是 ind-可表函子; 比较 \hyperref[I.cor8.13.3]{8.13.3}.
对函子 $F:C'\rightarrow C''$, 当 $C'$ 与 $C''$ 重新为任意 $𝒰$-范畴时,
我们将在下文 \hyperref[I.sec8.11.5]{(8.11.5)} 中更一般地提出一个
“改进”左正合性的概念 (即容许 pro-伴随的函子的概念).
\end{enonce*}

% END INCLUDED SGA4_English_current_buildable_source_20260906/source/01/drafts/SGA4_I_LINES_3228_3408_English.texfrag


% BEGIN INCLUDED SGA4_English_current_buildable_source_20260906/source/01/drafts/SGA4_I_LINES_3409_3598_English.texfrag
\subsection{常值与本质常值 ind-对象}\etiquette[8.4]{I.sec8.4}

选取一个终范畴 $\sheaf{e}$, 即 $\Ob
% 2013-5-24 compilation problem with \sheaf \XXX
\sheaf{e}$ 和 $F\ell$ $\sheaf{e}$ 各仅含一个元素的范畴
(例如, 若希望作典范选择, 可取 $\sheaf{e}$ 为 $(\Ens)$ 中以空集为唯一对象的满子范畴).
这显然是小滤过范畴. 若 $X$ 为 $C$ 的对象, 则将以 $\sheaf{e}$ 为指标、
取值为 $X$ 的常值 ind-对象与之对应, 记为 $c(X)$.
显然, 当 $X$ 变动时, 这给出全忠实函子
\hypertarget{sga4.exp.I.equation.8.4.1}{}$$
\displaylines{\rlap{$(8.4.1)$}\hfill c:C\longrightarrow
\Ind(C){\quoi},\hfill}
$$
且它在对象上单射, 我们借助它将 $C$ 等同于 $\Ind(C)$ 的满子范畴.
注意复合函子
\hypertarget{sga4.exp.I.equation.8.4.2}{}$$
\displaylines{\rlap{$(8.4.2)$}\hfill
C\xrightarrow{c}\Ind(C)\xrightarrow{L}\widehat{C}\hfill}
$$
就是典范函子 $h$ \hyperlink{sga4.exp.I.equation.8.2.1.1}{(8.2.1.1)}.

\setcounter{subsubsection}{2}
\subsubsection{}\etiquette[8.4.3]{I.sec8.4.3}
更\pageoriginale{}一般地, 设 $I$ 为滤过范畴.
$I$ 上取值为 $C$ 的对象 $X$ 的常值函子,
也称为\emph{以 $I$ 为指标、取值为 $X$ 的常值 ind-对象}.
显然, 若 $I$ 本质小, 此 ind-对象 ind-表示由 $X$ 表示的函子 $h(X)$,
故此 ind-对象同构于 $c(X)$.

\subsubsection{}\etiquette[8.4.4]{I.sec8.4.4}
如果 $C$ 的一个 ind-对象 $𝒳$ 同构于形如 $c(X)$, $X\in \ob C$ 的 ind-对象
(在适当的 $𝒱$, $𝒱'$ 所给的范畴 $\Ind_{𝒱}(C,𝒱')$ 中),
就称它 \emph{本质常值}.
此时对象 $X$ 确定到典范同构, 称为所讨论的本质常值 ind-对象的\emph{值}.
因此, 按定义, \hyperlink{sga4.exp.I.equation.8.4.1}{(8.4.1)} 的函子 $c$
诱导 $C$ 与 $\Ind(C)$ 中由本质常值 ind-对象组成的满子范畴的等价;
该子范畴即函子 $c$ 的本质像.

显然, 常值 ind-对象本质常值, 而逆命题仅在 $C$ 为空或为单点范畴这一平凡情形中成立.

\subsection{$\Ind(C)$ 中的滤过归纳极限}\etiquette[8.5]{I.sec8.5}

\begin{enonce*}{命题 8.5.1}\etiquette[8.5.1]{I.prop8.5.1}
设 $C$ 为 $𝒰$-范畴. 在 $\Ind(C)$ 中, 小滤过归纳极限可表,
且典范函子 $\hyperlink{sga4.exp.I.equation.8.2.4.8}{(8.2.4.8)}$
$$
L:\Ind(C)\longrightarrow \widehat{C}
$$
与它们交换.
\end{enonce*}

由于\pageoriginale $L$ 全忠实, 命题中的断言等价于下述断言:

\begin{enonce*}{推论 8.5.2}\etiquette[8.5.2]{I.cor8.5.2}
ind-可表预层的任意小滤过归纳极限 (在 $\widehat{C}$ 中取得) 都 ind-可表.
\end{enonce*}

这容易由判别准则 \hyperref[I.them8.3.3]{8.3.3} (ii) 得出.
验证的细节留待读者完成.

\setcounter{subsubsection}{2}
\subsubsection{}\etiquette[8.5.3]{I.sec8.5.3}
$\Ind(C)$ 中滤过归纳极限的“重言式”情形是从 $C$ 的一个本质小 ind-对象
$𝒳=(X_i)_{i\in I}$ 出发. 此时在 $\widehat{C}$ 中,
从而也在 $\Ind(C)$ 中 (或在 $\widehat{C}$ 中), 有:
\hypertarget{sga4.exp.I.equation.8.5.3.1}{}$$
\displaylines{\rlap{$(8.5.3.1)$}\hfill 𝒳=\varinjlim_I
X_i{\quoi}.\hfill}
$$
写此公式时应注意, 这不是 $C$ 中的归纳极限;
甚至当后者存在时, 它在 $\Ind(C)$ 中也不同构于 $𝒳$:
事实上, 典范函子 \hyperlink{sga4.exp.I.equation.8.4.1}{(8.4.1)}
$c:C\rightarrow \Ind(C)$ \emph{一般不与滤过归纳极限交换.}
(参见下文 \hyperref[I.exer8.5.5]{8.5.5}.)

为避免写 \hyperlink{sga4.exp.I.equation.8.5.3.1}{(8.5.3.1)} 时可能产生的这种混淆,
“某些作者” (= {P. Deligne}) 倾向于将它写作
\hypertarget{sga4.exp.I.equation.8.5.3.2}{}$$
\displaylines{\rlap{$(8.5.3.2)$}\hfill
%𝒳=\lq\lq{\varinjlim_I}\textrm{ »\,}
𝒳=\textrm{« }{\varinjlim_I}\textrm{ »\,}X_i{\quoi},\hfill}
$$
引号表示归纳极限在 ind-对象范畴 $\Ind(C)$ 中取得.
进而, $\Ind(C)$ 中的每个归纳极限运算都应记为
$\textrm{« }{\lim}\textrm{ »}$
(所考虑的 $\Ind(C)$ 中的归纳系的分量不一定属于
$c:C\rightarrow \Ind(C)$ 的像或本质像).

\subsubsection{}\etiquette[8.5.4]{I.sec8.5.4}
为\pageoriginale{}计算函子
\hypertarget{sga4.exp.I.equation.8.5.4.1}{}$$
\displaylines{\rlap{$(8.5.4.1)$}\hfill \varphi:J\longrightarrow
\Ind(C){\quoi},\hfill}
$$
的归纳极限或投射极限 (这里 $J$ 为任意范畴), 常常便于用函子
\hypertarget{sga4.exp.I.equation.8.5.4.2}{}$$
\displaylines{\rlap{$(8.5.4.2)$}\hfill \Phi:J\times I\longrightarrow
C{\quoi},\hfill}
$$
将 $\varphi$ 明确表示出来, 其中 $I$ 为本质小滤过范畴, 最好为小范畴.
这样的 $\Phi$ 确实定义函子
$$
J\longmapsto (i\longmapsto \Phi (j,i)):J\longmapsto
\Hhom(I,C){\quoi},
$$
从而与典范箭头 \hyperlink{sga4.exp.I.equation.8.2.6.1}{(8.2.6.1)}
$\Hhom(I,C)\rightarrow \Ind(C)$ 复合, 得到函子
\hypertarget{sga4.exp.I.equation.8.5.4.3}{}$$
\displaylines{\rlap{$(8.5.4.3)$}\hfill J\longmapsto (i\longmapsto
\Phi(j,i)):J\longrightarrow \Ind(C){\quoi}.\hfill}
$$
如果对应的函子 \hyperlink{sga4.exp.I.equation.8.5.4.3}{(8.5.4.3)} 同构于 $\varphi$,
就称 $\Phi$ 为\emph{以 (滤过) 范畴 $I$ 为指标的 $\varphi$ 的带指标表达式}.
下文 (\hyperref[I.sec8.8]{8.8}) 将研究给定函子 $\varphi$ 存在带指标表达式的一般条件;
这里则限于从以带指标形式 \hyperlink{sga4.exp.I.equation.8.5.4.3}{(8.5.4.3)} 给出的函子
$\varphi$ 出发, 并假定 $J$ 为小滤过范畴,
以说明在此情形中如何计算
$\textrm{«}{\varinjlim}\textrm{»\,}\varphi=\textrm{«}{\varinjlim_j}\textrm{»\,}\varphi(j)$.
公式 \hyperlink{sga4.exp.I.equation.8.5.3.2}{(8.5.3.2)} 连同归纳极限的结合律公式
(\hyperref[I.sec2.5.0]{2.5.0}) 立即给出
$\textrm{«}{\varinjlim}\textrm{»\,}\varphi$ 的“\emph{重言式计算}”:
\hypertarget{sga4.exp.I.equation.8.5.4.4}{}$$
\displaylines{\rlap{$(8.5.4.4)$}\hfill
\textrm{«}{\varinjlim_J}\textrm{»\,}\varphi=\textrm{«}{\varinjlim_{J\times
I}}\textrm{»\,}\Phi{\quoi},\hfill}
$$
即\pageoriginale
\hypertarget{sga4.exp.I.equation.8.5.4.5}{}$$
\displaylines{\rlap{$(8.5.4.5)$}\hfill
\textrm{«}{\varinjlim_j}\textrm{»\,}(i\longmapsto \Phi
(j,i))=\textrm{«}{\varinjlim_{j,i}}\textrm{»\,}\Phi(j,i){\quoi},\hfill}
$$
因此, 所求的归纳系正是 $\Phi$ 本身
(指标范畴为 $J\times I$, 它确实滤过, 因为 $J$ 与 $I$ 均滤过).

\begin{enonce*}[remark]{习题 8.5.5}\etiquette[8.5.5]{I.exer8.5.5}
设 $C$ 为 $𝒰$-范畴.
\begin{enumerate}
\renewcommand{\theenumi}{\alph{enumi}}
\renewcommand{\labelenumi}{\theenumi)}
\item 证明下述条件等价:

\begin{itemize}
\item[(i)] 在 $C$ 中, 小归纳极限可表, 且函子 $c:C\rightarrow \Ind(C)$ 与它们交换.

\item[(iii)] 函子 $c$ 为范畴等价.

\item[(ii)] 在 $C$ 中, 小归纳极限可表, 且对每个 $X\in \Ob C$,
函子 $Y\mapsto \Hom(X,Y)$ 与这些极限交换.

\item[(iv)] 小滤过归纳极限在 $C$ 中可表, 且函子 $\lim c:\Ind C\rightarrow C$
(参见 \hyperlink{sga4.exp.I.equation.8.7.1.5}{8.7.1.5}) 全忠实 (等价地, 为等价).
\end{itemize}

\item 若 $C$ 为有限范畴, 证明上述条件等价于:
对 $C$ 中的每个投影算子 (\hyperref[I.sec10.6]{10.6}), 其像在 $C$ 中可表.
\end{enumerate}
\end{enonce*}

% END INCLUDED SGA4_English_current_buildable_source_20260906/source/01/drafts/SGA4_I_LINES_3409_3598_English.texfrag


% BEGIN INCLUDED SGA4_English_current_buildable_source_20260906/source/01/drafts/SGA4_I_LINES_3599_3796_English.texfrag
\subsection{函子向 ind-对象的扩张}\etiquette[8.6]{I.sec8.6}

\subsubsection{}\etiquette[8.6.1]{I.sec8.6.1}
设
\hypertarget{sga4.exp.I.equation.8.6.1.1}{}$$
\displaylines{\rlap{$(8.6.1.1)$}\hfill f:C\longrightarrow C'\hfill}
$$
为 $𝒰$-范畴之间的函子. 对 $C$ 的每个 ind-对象
$$
\varphi:I\longrightarrow C
$$
(\pageoriginale $I$ 为滤过范畴), 复合函子 $f\circ\varphi$ 是 $C'$ 的 ind-对象,
也记为 $\Ind(f)(\varphi)$. 用带指标的记号, 若将 $\varphi$ 写为
$𝒳=(X_i)_{i\in I}$, 则得到
\hypertarget{sga4.exp.I.equation.8.6.1.2}{}$$
\displaylines{\rlap{$(8.6.1.2)$}\hfill \Ind(f)(X_i)_{i\in
I}=(f(X_i))_{i\in I}{\quoi}.\hfill}
$$
设 $𝒴=(Y_j)_{j\in J}$ 为 $C$ 的另一个归纳系.
于是得到显然的映射, 也记为 $u\rightarrow\Ind(f)(u)$:
\hypertarget{sga4.exp.I.equation.8.6.1.3}{}$$
\displaylines{\rlap{$(8.6.1.3)$}\hfill
\Hom(𝒳,𝒴)=\varprojlim_i \varinjlim_j \Hom
(X_i,Y_j)\rightarrow\hfill\cr \hfill
\Hom(\Ind(f)(𝒳),\Ind(f)(𝒴))\simeq \varprojlim_i
\varinjlim_j \Hom(f(X_i),f(Y_j)){\quoi}.\hfill}
$$
参照未写出的 ind-对象态射的复合公式 (\hyperref[I.sec8.2.5]{8.2.5}),
立即可见这些映射与 ind-对象态射的复合相容.
因而, 对每个宇宙 $𝒱$, 我们得到了函子
\hypertarget{sga4.exp.I.equation.8.6.1.4}{}$$
\displaylines{\rlap{$(8.6.1.4)$}\hfill
\Ind(f):\Ind_{𝒱}(C,𝒰)\longrightarrow
\Ind_{𝒱}(C',𝒰)\hfill}
$$
(记号见 \hyperlink{sga4.exp.I.equation.8.2.4.5}{(8.2.4.5)}), 特别得到函子
\hypertarget{sga4.exp.I.equation.8.6.1.5}{}$$
\displaylines{\rlap{$(8.6.1.5)$}\hfill
\Ind(f):\Ind(C)\longrightarrow \Ind(C'){\quoi}.\hfill}
$$
若还有另一个函子
$$
g:C'\longrightarrow C''{\quoi},
$$
则显然有如下 $\Ind$-范畴 (或按需取 $\Ind_{𝒱}$-范畴) 之间的函子恒等式:
\hypertarget{sga4.exp.I.equation.8.6.1.6}{}$$
\displaylines{\rlap{$(8.6.1.6)$}\hfill \Ind(gf)=\Ind(g)\circ
\Ind(f){\quoi},\hfill}
$$
并且,\pageoriginale{}当 $C'=C$ 时, 有一贯的公式
\hypertarget{sga4.exp.I.equation.8.6.1.7}{}$$
\displaylines{\rlap{$(8.6.1.7)$}\hfill
\Ind(\id_C)=\id_{\Ind(C)}{\quoi}.\hfill}
$$
因而, 启发式地说, 范畴 $\Ind(C)$ \emph{函子性地依赖于 $C$} (协变地)
\footnote[1]{关于某些条件下的反变依赖性, 参见 \hyperref[I.sec8.11]{8.11}.}.
请读者通过对每对 $𝒰$-范畴 $C$, $C'$ 定义典范函子
\hypertarget{sga4.exp.I.equation.8.6.1.8}{}$$
\displaylines{\rlap{$(8.6.1.8)$}\hfill
\Iind(f):\Hhom(C,C')\longrightarrow \Hhom(\Ind(C),
\Ind(C')){\quoi},\hfill}
$$
并指明恒等同构 \hyperlink{sga4.exp.I.equation.8.6.1.6}{(8.6.1.6)}
与 \hyperlink{sga4.exp.I.equation.8.6.1.7}{(8.6.1.7)} 的函子性,
将其进一步明确为 2-\emph{函子性}的依赖关系.

\subsubsection{}\etiquette[8.6.2]{I.sec8.6.2}
回到函子 \hyperlink{sga4.exp.I.equation.8.6.1.1}{(8.6.1.1)}, 考虑相应的函子
\hyperref[I.prop5.1]{(5.1)}
\hypertarget{sga4.exp.I.equation.8.6.2.1}{}$$
\displaylines{\rlap{$(8.6.2.1)$}\hfill
f_{!}:\widehat{C}\longrightarrow \widehat{C'}{\quoi},\hfill}
$$
以及函子图
\hypertarget{sga4.exp.I.equation.8.6.2.2}{}\[\vcenter{
\xymatrix@C=3.25cm@R=1.25cm{
\Ind(C) \ar[r]^{\ind(f)} \ar[d]_{L_C}& \Ind(C')\ar[d]^{L_{C'}}\\
\hat{C} \ar[r]^{f_!}& \hat{C}' }}\leqno{(8.6.2.2)}
\]
其中竖箭头表示 \hyperlink{sga4.exp.I.equation.8.2.4.8}{(8.2.4.8)}
的典范函子 $L$. 由于函子 $f_{!}$ “扩张 $f$”并与归纳极限交换
\hyperlink{sga4.exp.I.proposition.5.4.item.3}{(5.4.3)},
而函子 $L_C$ 和 $L_{C'}$ 与小滤过归纳极限交换
(\hyperref[I.prop8.5.1]{8.5.1}), 由
\hyperlink{sga4.exp.I.equation.8.5.3.2}{(8.5.3.2)} 得出,
对 $C$ 的每个 ind-对象 $(X_i)_{i\in I}$, 在 $\widehat{C'}$ 中有典范同构
$$
f_{!}(L_C((X_i)_{i\in I}))\simeq \varinjlim_i f_{!}L_C(X_i)\simeq
\varinjlim_i L_{C'}f(X_i)=L_{C'}(f(X_i)_{i\in I})
$$
即\pageoriginale{}典范同构
$$
f_{!}L_C(𝒳)\simeq L_{C'}(\Ind(f)(𝒳)){\quoi}.
$$
此同构的函子性留待读者验证, 即它对应于典范同构
\hypertarget{sga4.exp.I.equation.8.6.2.3}{}$$
\displaylines{\rlap{$(8.6.2.3)$}\hfill f_{!}L_C\simeq L_{C'}\circ
\Ind(f){\quoi}.\hfill}
$$
换言之, \emph{方块 \hyperlink{sga4.exp.I.equation.8.6.2.2}{(8.6.2.2)}
交换, 精确到典范同构.} 考虑到 $f_{!}$ 与小归纳极限交换,
并利用 \hyperref[I.prop8.5.1]{8.5.1}, 得到:

\begin{enonce*}{命题 8.6.3}\etiquette[8.6.3]{I.prop8.6.3}
设 $f:C\rightarrow C'$ 为 $𝒰$-范畴之间的函子. 则函子
$$
\Ind(f):\Ind(C)\longrightarrow \Ind(C')
$$
与滤过归纳极限交换. 此外, 它使下图交换:
\[
\xymatrix@C=2.5cm@R=1.5cm{
C       \ar[r]^{f}       \ar[d]_{c_{C}}& C' \ar[d]^{c_{C'}}\\
\Ind(C) \ar[r]^{\Ind(f)} & \Ind(C')\quad , }
\]
其中竖箭头 $c_C$, $c_{C'}$ 为典范函子
$\hyperlink{sga4.exp.I.equation.8.4.1}{(8.4.1)}$.
\end{enonce*}
最后一个断言由定义显然成立, 列出它只是为了便于引用.
还应注意, \hyperref[I.prop8.6.3]{8.6.3} 所述的性质\emph{刻画了函子}
$\Ind(f)$, 精确到唯一同构,\pageoriginale{}即作为
$f_{!}$ \hyperlink{sga4.exp.I.equation.8.6.2.1}{(8.6.2.1)} 诱导的函子;
这由刚才给出的 \hyperlink{sga4.exp.I.equation.8.6.2.3}{(8.6.2.3)} 的证明得出.

\begin{enonce*}{命题 8.6.4}\etiquette[8.6.4]{I.prop8.6.4}
记号同 $\hyperref[I.prop8.6.3]{8.6.3}$.
\begin{enumerate}
\renewcommand{\theenumi}{\alph{enumi}}
\renewcommand{\labelenumi}{\theenumi)}
\item 函子 $\Ind(f)$ 忠实 (分别地, 全忠实) 当且仅当 $f$ 忠实 (分别地, 全忠实).

\item 函子 $\Ind(f)$ 为范畴等价, 当且仅当 $f$ 全忠实,
且 $C'$ 的每个对象同构于 $f$ 的像中某个对象的直因子 \hyperref[I.sec10.6]{(10.6)}.
\end{enumerate}
\end{enonce*}

\noindent \emph{证明.}
\begin{enumerate}
\renewcommand{\theenumi}{\alph{enumi}}
\renewcommand{\labelenumi}{\theenumi)}
\item 必要性显然来自 $f$ 由 $\Ind(f)$ 诱导这一事实.
对充分性, 只须使用 $\Ind(f)$ 在集合 $\Hom(𝒳,𝒴)$ 上的形式
\hyperlink{sga4.exp.I.equation.8.6.1.3}{(8.6.1.3)}, 并回顾集合的滤过归纳极限
以及任意投射极限均将单态射映为单态射、将同构映为同构.

\item 已可假定 $f$, 从而 $\Ind(f)$, 全忠实.
由于 $\Ind(C')$ 的每个对象都是 $C'$ 的对象的小滤过归纳极限,
$\Ind(f)$ 的全忠实性表明, 此函子本质满射等价于 $C'$ 的每个对象 $X'$ 都属于其本质像.
但若有同构
$$
X'\xrightarrow{\sim} \textrm{« }{\varinjlim}\textrm{ »\,}f(X_i){\quoi},
$$
此同构经某个 $f(X_i)$ 分解, 表明 $X'$ 是 $f(X_i)$ 的直因子,
由此证明了 b) 的必要性. 对充分性, 利用 $\Ind(f)$ 的全忠实性, 断言立即成立.
\end{enumerate}

\begin{enonce*}{推论 8.6.5}\etiquette[8.6.5]{I.cor8.6.5}
在\pageoriginale $\Ind(E)$ 中, 投影算子
(\hyperref[I.sec10.6]{10.6}) 的像可表,
且函子 $\Ind(f):\Ind(C)\rightarrow \Ind(C')$ 与这些像的形成交换.
\end{enonce*}

事实上, 这来自小滤过归纳极限在 $\Ind(C)$ 中可表
(\hyperref[I.prop8.5.1]{8.5.1}), 且 $\Ind(f)$ 与它们交换
(\hyperref[I.prop8.6.3]{8.6.3}) 这两点, 并须注意投影算子
$p:X\rightarrow X$ 的像可解释为如下以 $\NN$ 为指标的滤过归纳系的归纳极限
$$
X\xrightarrow{P}X\xrightarrow{P}X\longrightarrow\ldots
$$
或者等价地, 解释为滤过范畴 $P$ 上的显然函子的归纳极限,
其中只有一个对象及一条非恒等箭头 $\Pi$, 满足 $\Pi^2=\Pi$.


% END INCLUDED SGA4_English_current_buildable_source_20260906/source/01/drafts/SGA4_I_LINES_3599_3796_English.texfrag


% BEGIN INCLUDED SGA4_English_current_buildable_source_20260906/source/01/drafts/SGA4_I_LINES_3797_3955_English.texfrag
\subsection{}\etiquette[8.7]{I.sec8.7}
\emph{函子 $\varinjlim_C:\Ind(C)\rightarrow C$.
范畴 $\Ind(C)$ 的泛性质刻画.}

\subsubsection{}\etiquette[8.7.1]{I.sec8.7.1}
回到 $𝒰$-范畴 $C$ 以及典范函子
\hypertarget{sga4.exp.I.equation.8.7.1.1}{}$$
\displaylines{\rlap{$(8.7.1.1)$}\hfill c:C\longrightarrow
\Ind(C){\quoi}.\hfill}
$$
设 $𝒳=(X_i)_{i\in I}$ 为 $\Ind(C)$ 的对象.
由可表归纳极限的定义 (\hyperref[I.def2.1]{2.1}, \hyperref[I.sec2.1.1]{2.1.1})
立即得出: \emph{$\varinjlim_iX_i$ 在 $C$ 中可表, 当且仅当 $c$ 的左伴随在 $𝒳$ 处有定义},
且此时 $\varinjlim_iX_i$ (在 $C$ 中计算) 正是该左伴随在 $𝒳$ 处的值:
\hypertarget{sga4.exp.I.equation.8.7.1.2}{}$$
\displaylines{\rlap{$(8.7.1.2)$}\hfill \Hom_C(\varinjlim_i X_i,
Y)\simeq \Hom_{\Ind(C)}((X_i)_{i\in I}, c(Y)){\quoi}.\hfill}
$$
若\pageoriginale{}用
\hypertarget{sga4.exp.I.equation.8.7.1.3}{}$$
\displaylines{\rlap{$(8.7.1.3)$}\hfill \Ind(C)'\subset
\Ind(C)\hfill}
$$
表示 $\Ind(C)$ 中由那些\emph{在} $C$ \emph{中}有归纳极限的 $C$ 的 ind-对象组成的满子范畴,
则上述观察表明此子范畴\emph{严格满}, 并且存在典范函子
\hypertarget{sga4.exp.I.equation.8.7.1.4}{}$$
\displaylines{\rlap{$(8.7.1.4)$}\hfill
\textstyle{\varinjlim}_C:\Ind(C)'\longrightarrow C{\quoi},\hfill}
$$
其在 $\Ind(C)'$ 的每个对象 $(X_i)_{i\in I}$ 处的值为该对象在 $C$ 中的归纳极限.
特别地, 当小滤过归纳极限在 $C$ 中可表时, 就得到自然函子
\hypertarget{sga4.exp.I.equation.8.7.1.5}{}$$
\displaylines{\rlap{$(8.7.1.5)$}\hfill
{\varinjlim}_C:\Ind(C)\longrightarrow C{\quoi}.\hfill}
$$

\setcounter{paragraph}{5}
\paragraph{}\etiquette[8.7.1.6]{I.sec8.7.1.6}
当然, 上述函子也可定义在形如 $\Ind_{𝒱}(C,𝒰)'$ 和 $\Ind_{𝒱}(C,𝒰)$ 的范畴上,
但由 $\hyperlink{sga4.exp.I.equation.8.2.4.6}{(8.2.4.6)}$,
它们已由上述对应于 $𝒱=𝒰$ 情形的函子确定 (精确到唯一同构).

\paragraph{}\etiquette[8.7.1.7]{I.sec8.7.1.7}
上述将 $\varinjlim_C$ 构造为左伴随函子的方法立即表明,
此函子\emph{与任意归纳极限交换}, 特别是与小滤过归纳极限交换
(后者在 $\Ind(C)$ 中可表 (\hyperref[I.prop8.5.1]{8.5.1})).
还应注意, $\Ind(C)'$ 总包含 $c$ 的本质像 (由本质常值 ind-对象组成),
且存在对 $X\in\ob\Ind(C)'$ 函子性的典范同构:
\hypertarget{sga4.exp.I.equation.8.7.1.8}{}$$
\displaylines{\rlap{$(8.7.1.8)$}\hfill {\varinjlim}_C c(X)\simeq
X{\quoi},\hfill}
$$
即,\pageoriginale{}在 $C$ 对小归纳极限稳定这一有利情形中, 有典范同构
\hypertarget{sga4.exp.I.equation.8.7.1.9}{}$$
\displaylines{\rlap{$(8.7.1.9)$}\hfill {\varinjlim}_C \circ c\simeq
\id_C{\quoi}.\hfill}
$$

\subsubsection{}\etiquette[8.7.2]{I.sec8.7.2}

现考虑函子
\hypertarget{sga4.exp.I.equation.8.7.2.1}{}$$
\displaylines{\rlap{$(8.7.2.1)$}\hfill f:C\longrightarrow
E{\quoi},\hfill}
$$
其中 $E$ 为小归纳极限可表的 $𝒰$-范畴, 从而存在函子
$$
{\varinjlim}_E:\Ind(E)\longrightarrow E{\quoi}.
$$
由于我们也已定义 \hyperref[I.sec8.6]{(8.6)}
$$
\Ind(f):\Ind(C)\longrightarrow \Ind(E){\quoi},
$$
可考虑复合
\hypertarget{sga4.exp.I.equation.8.7.2.2}{}$$
\displaylines{\rlap{$(8.7.2.2)$}\hfill \bar{f}={\varinjlim}_E\circ
\Ind(f):\Ind(C)\longrightarrow E{\quoi},\hfill}
$$
有时称为 \emph{$f$ 向 ind-对象的典范扩张}
(但不可与 $\Ind(f)$ 混淆!). 它是两个与小滤过归纳极限交换的函子
(\hyperref[I.prop8.6.3]{8.6.3}, \hyperref[I.sec8.7.1.7]{8.7.1.7}) 的复合,
因而本身也与小滤过归纳极限交换.
此外, 将同构 \hyperlink{sga4.exp.I.equation.8.7.1.9}{(8.7.1.9)} 应用于 $E$,
连同 \hyperlink{sga4.exp.I.equation.8.6.2.3}{(8.6.2.3)},
表明 $\bar f$ \emph{扩张} $f$, 精确到同构, 即存在典范同构
\hypertarget{sga4.exp.I.equation.8.7.2.3}{}$$
\displaylines{\rlap{$(8.7.2.3)$}\hfill \bar{f}\circ c_C\simeq
f{\quoi}.\hfill}
$$
又由于\pageoriginale $\Ind(C)$ 的每个对象都是 $C$ 的对象的小滤过归纳极限,
颇易看出, 上述两个性质仍然\emph{刻画} $\bar f$, 精确到唯一同构.
将此结论加强, 便可同时在小滤过归纳极限可表的 $𝒰$-范畴 $E$ 中,
得到 $\Ind(C)$ 的泛性质刻画, 精确到等价.
若 $E$ 与 $F$ 为两个这样的 $𝒰$-范畴, 用
\hypertarget{sga4.exp.I.equation.8.7.2.4}{}$$
\displaylines{\rlap{$(8.7.2.4)$}\hfill \Hhom(E,F)'\subset
\Hhom(E,F)\hfill}
$$
表示 $\Hhom(E,F)$ 中由与小滤过归纳极限交换的函子组成的满子范畴. 则有:

\begin{enonce*}{命题 8.7.3}\etiquette[8.7.3]{I.prop8.7.3}
设 $C$ 为 $𝒰$-范畴, 采用上面的记号
\hyperlink{sga4.exp.I.equation.8.7.2.4}{(8.7.2.4)}. 则典范函子
$$
c_C:C\longrightarrow \Ind(C)
$$
在以 $C$ 为起点、以小滤过归纳极限可表的 $𝒰$-范畴为终点的函子中为 $2$-泛的.
换言之, $\Ind(C)$ 是这样的范畴
(\hyperref[I.sec8.2.5]{8.2.5}, \hyperref[I.prop8.5.1]{8.5.1}),
且若 $E$ 为这样的范畴, 则函子
\hypertarget{sga4.exp.I.equation.8.7.3.1}{}$$
\displaylines{\rlap{$(8.7.3.1)$}\hfill g\longmapsto g\circ
c_C:\Hhom(\Ind(C),E)'\longrightarrow \Hhom(C,E)\hfill}
$$
为范畴等价.
\end{enonce*}

为证最后一点, 只须验证由 \hyperlink{sga4.exp.I.equation.8.7.2.2}{(8.7.2.2)}
定义的映射 $f\mapsto\bar f$ 可细化为\emph{函子}
\hypertarget{sga4.exp.I.equation.8.7.3.2}{}$$
\displaylines{\rlap{$(8.7.3.2)$}\hfill f\longmapsto
\bar{f}={\varinjlim}_E\cdot \Ind(f):\Hhom(C,E)\longrightarrow
\Hhom(\Ind(C),E)'{\quoi},\hfill}
$$
并且\pageoriginale{}它是 \hyperlink{sga4.exp.I.equation.8.7.3.1}{(8.7.3.1)} 的拟逆.
这一实质上平凡的验证的细节留待读者完成.
也可先用 \hyperref[I.prop7.8]{7.8} 得出
\hyperlink{sga4.exp.I.equation.8.7.3.1}{(8.7.3.1)} 全忠实,
再用 \hyperlink{sga4.exp.I.equation.8.7.2.3}{(8.7.2.3)} 得出它本质满射.

\setcounter{subsubsection}{3}

% END INCLUDED SGA4_English_current_buildable_source_20260906/source/01/drafts/SGA4_I_LINES_3797_3955_English.texfrag


% BEGIN INCLUDED SGA4_English_current_buildable_source_20260906/source/01/drafts/SGA4_I_LINES_3956_4068_English.texfrag
\subsubsection{}\etiquette[8.7.4]{I.sec8.7.4}
回到 $𝒰$-范畴之间的函子
\hypertarget{sga4.exp.I.equation.8.7.4.1}{}$$
\displaylines{\rlap{$(8.7.4.1)$}\hfill f:C\longrightarrow
E{\quoi},\hfill}
$$
其中 $E$ 为小滤过归纳极限可表的 $𝒰$-范畴, 由此得到函子
\hypertarget{sga4.exp.I.equation.8.7.4.2}{}$$
\displaylines{\rlap{$(8.7.4.2)$}\hfill \bar{f}:(X_i)_{i\in I}
\longmapsto \mathop{\lim}_i f(X_i):\Ind(C)\longrightarrow
E{\quoi}.\hfill}
$$
我们要研究保证 $\bar f$ 全忠实或分别为范畴等价的 $f$ 的条件.
由于 $f$ 同构于复合 $\bar f\circ c_C$, 且 $c_C:C\rightarrow\Ind(C)$ 全忠实,
可见 $\bar f$ 全忠实的\emph{必要}条件是 $f$ 全忠实.
因此, 用 $C$ 在 $E$ 中的本质像替换之, 就可设 $f$ 为 $E$ 的某个子范畴 $C$ 的包含函子,
而实质上不失一般性; 下文将如此假定以简化记号.

\begin{enonce*}{命题 8.7.5}\etiquette[8.7.5]{I.prop8.7.5}
记号同 \hyperref[I.sec8.7.4]{8.7.4}.
\begin{enumerate}
\renewcommand{\theenumi}{\alph{enumi}}
\renewcommand{\labelenumi}{\theenumi)}
\item 函子 $\bar f$ 全忠实的充要条件是:
\begin{itemize}
\item[(i)] $C$ 的每个对象 $X$ 均满足下述条件:

{\rm(PF)} 对 $C$ 中每个归纳极限 $\varinjlim_iY_i$ 在\pageoriginale $E$ 中存在的
小滤过归纳系 $(Y_i)_{i\in I}$, 典范映射
\hypertarget{sga4.exp.I.equation.8.7.5.1}{}$$
\displaylines{\rlap{$(8.7.5.1)$}\hfill \varinjlim_i \Hom
(X,Y_i)\longrightarrow \Hom(X,\varinjlim_i Y_i)\hfill}
$$
为双射.
\end{itemize}

\item 函子 $\bar f$ 为范畴等价, 当且仅当 $C$ 满足条件 {\rm(i)} 及下述两个条件:
\begin{itemize}
\item[(ii)] $C$ 是 $E$ 的由严格满态射生成的子范畴 (\hyperref[I.def7.1]{7.1}).

\item[\rm(iii)] 对 $E$ 的每个对象 $X$, $E_{/X}$ 中由那些起点属于 $\Ob C$ 的
$X$ 上对象 $X'$ 组成的满子范畴 $C_{/X}$ 滤过且本质小.
\end{itemize}
\end{enumerate}

特别地, 若 $C$ 等价于小范畴, $C$ 中的有限归纳极限可表,
且包含函子 $f:C\rightarrow E$ 与它们交换
(即对每个有限范畴 $J$ 和每个函子 $\varphi:J\rightarrow C$,
$f\varphi$ 的归纳极限可表, 且在 $E$ 中同构于 $C$ 的对象), 则条件 {\rm(iii)} 成立.
\end{enonce*}


\noindent \emph{证明}.
\begin{enumerate}
\renewcommand{\theenumi}{\alph{enumi}}
\renewcommand{\labelenumi}{\theenumi)}
\item 采用条件 (i) 的记号, 用 $𝒴$ 表示 ind-对象 $(Y_i)$,
用 $𝒳$ 表示由 $X$ 定义的常值 ind-对象.
则 \hyperlink{sga4.exp.I.equation.8.7.5.1}{(8.7.5.1)} 正是典范映射
$$
\Hom(𝒳, 𝒴)\longrightarrow \Hom
(\bar{f}(𝒳),\bar{f}(𝒴)){\quoi},
$$
所以若 $\bar f$ 全忠实, 此映射确实为双射.
反之, 若 $𝒳=(X_j)_{j\in J}$ 与 $𝒴=(Y_i)_{i\in I}$ 是 $E$ 的任意 ind-对象,
则映射 $u:\Hom(𝒳,𝒴)\rightarrow\Hom(\bar f(𝒳),\bar f(𝒴))$
是\pageoriginale{}各映射
$u_j:\Hom(𝒳_j,𝒴)\rightarrow\Hom(\bar f(𝒳_j),𝒴)$ 对 $j$ 的投射极限,
其中 $𝒳_j$ 为取值为 $X_j$ 的常值 ind-对象.
因此, 若各 $u_j$ 为双射, $u$ 便为双射, 从而 (i) 蕴含 $\bar f$ 全忠实.

\item 假定 (i)、(ii) 和 (iii) 成立, 证明 $\bar f$ 为等价.
只须证它本质满射, 即每个 $X\in\ob E$ 属于其本质像.
然而假设 (ii) 表示 $\varinjlim_{C_{/X}}Z\rightarrow X$ 是同构,
而 (iii) 表示范畴 $C_{/X}$ 滤过且本质小.
故 $X$ 是由自然函子 $C_{/X}\rightarrow E$ 定义的 $E$ 的 ind-对象的像.
反之, 设 $\bar f$ 为等价. 由 a), 只须验证 (ii) 与 (iii).
可设 $E=\Ind(C)$, 并将 $C$ 等同于 $E$ 的子范畴 $c(C)$.
我们知道 (\hyperref[I.them8.3.3]{8.3.3} (ii)), 对每个 $X\in\Ind(C)$,
如有需要, 将其等同于其所 ind-表示的 $C^\circ$ 上函子 $F$,
则 $C_{/X}=C_{/F}$ 为本质小滤过范畴 (这证明了 (iii)),
且函子 $C_{/F}\rightarrow\widehat{C}$ 在 $\widehat{C}$ 中的归纳极限为 $F$.
由于 $F=X$ 属于 $E$, 它在 $\widehat C$ 的满子范畴 $E$ 中也如此,
这就证明了 (ii). 余下须证 b) 的最后一个断言.
但关于 $C$ 的假设显然蕴含有限归纳极限在范畴 $C_{/X}$ 中可表;
\textit{从而}范畴 $C_{/X}$ 滤过.
\end{enumerate}


% END INCLUDED SGA4_English_current_buildable_source_20260906/source/01/drafts/SGA4_I_LINES_3956_4068_English.texfrag


% BEGIN INCLUDED SGA4_English_current_buildable_source_20260906/source/01/drafts/SGA4_I_LINES_4069_4270_English.texfrag
\begin{enonce*}[remark]{注 8.7.5.2}\etiquette[8.7.5.2]{I.rem8.7.5.2}
上述\pageoriginale{}论证更一般地表明, 当条件 (i) 满足时,
全忠实函子 $\bar f$ \hyperlink{sga4.exp.I.equation.8.7.4.2}{(8.7.4.2)}
的本质像由下述 $X\in\ob E$ 组成: 范畴 $C_{/X}$ 滤过且本质小
(若 $C$ 满足 b) 结尾所述的条件, 此条件自动满足),
\emph{并且}典范函子 $C_{/X}\rightarrow E$ 的归纳极限为 $X$.
\end{enonce*}

\begin{enonce*}{推论 8.7.6}\etiquette[8.7.6]{I.cor8.7.6}
令 $E_{PF}$ 为 $E$ 中由满足 \hyperref[I.prop8.7.5]{8.7.5} {\rm a)}
的条件 $(PF)$ 的对象组成的满子范畴.
则对每个 $J$ 为有限范畴、且在 $E$ 中有归纳极限 $X$ 的函子
$J\rightarrow E_{PF}$, 均有 $X\in\ob E_{PF}$.
从包含 $g:E_{PF}\hookrightarrow E$ 得出的典范函子
$(X_i)_{i\in I}\rightarrow\displaystyle{\mathop{\varinjlim}_i}{}_EX_i$
\hypertarget{sga4.exp.I.equation.8.7.6.1}{}$$
\displaylines{\rlap{$(8.7.6.1)$}\hfill
\bar{g}:\Ind(E_{PF})\longrightarrow E\hfill}
$$
全忠实. 若有限归纳极限在 $E$ 中可表且 $E_{PF}$ 等价于小范畴,
则函子 $\bar f$ 的本质像由 $E$ 的那些使典范态射
$\varinjlim_{(E_{PF})_{/X}}Y\rightarrow X$ 为同构的对象 $X$ 组成.
若进一步 $E$ 的子范畴 $E_{PF}$ 由严格满态射生成
(\hyperref[I.def7.1]{7.1}), 则函子
\hyperlink{sga4.exp.I.equation.8.7.6.1}{(8.7.6.1)} 为范畴等价.
\end{enonce*}

由 \hyperref[I.prop8.7.5]{8.7.5}, 并对第一个断言使用
\hyperref[I.prop2.8]{2.8}, 按上述排列次序, 所有断言都显然成立.

\begin{enonce*}{推论 8.7.7}\etiquette[8.7.7]{I.cor8.7.7}
假定\pageoriginale{}有限归纳极限在 $E$ 中可表.
设 $C$ 为 $E$ 的一个等价于小范畴的满子范畴, 考虑函子
$(X_i)_{i\in I}\rightarrow
{\displaystyle{\mathop{\varinjlim}_i}{}_E}X_i$:
\hypertarget{sga4.exp.I.equation.8.7.7.1}{}$$
\displaylines{\rlap{$(8.7.7.1)$}\hfill \bar{f}:\Ind(C)\rightarrow
E{\quoi}.\hfill}
$$
令 $E_{PF}$ 如 $\hyperref[I.cor8.7.6]{8.7.6}$ 中所定义. 下述条件等价:
\begin{itemize}
\item[(i)] 函子 $\bar f:\Ind(C)\rightarrow E$ 为等价.

\item[(ii)] 范畴 $C$ 在 $E$ 中由严格满态射生成, 包含于 $E_{PF}$,
且 $E_{PF}$ 的每个对象在 $E$ 中
(或等价地, 在 $E_{PF}$ 中 (\hyperref[I.cor8.7.6]{8.7.6}))
同构于 $C$ 的某个对象的直因子.

当 $E$ 的子范畴 $C$ 对直因子稳定 (\hyperref[I.sec10.6]{10.6}) 时,
上述条件还等价于下述条件:

\item[(ii bis)] $C=E_{PF}$, 且 $C$ 在 $E$ 中由严格满态射生成.

\item[(iii)] $E$ 的子范畴 $C$ 由严格满态射生成, 包含于 $E_{PF}$,
且有限归纳极限在其中可表.

若进一步有限归纳极限在 $E$ 中可表, 则这些条件还等价于

\item[(iii bis)] $E$ 的子范畴 $C$ 由严格满态射生成, 包含于 $E_{PF}$,
且在 $E$ 中对有限归纳极限稳定 (等价地, 对有限和与双箭头的余核稳定).
\end{itemize}
\end{enonce*}

我们\pageoriginale{}已知 (\hyperref[I.prop8.7.5]{8.7.5}), 条件 (i) 蕴含
$C\subset E_{PF}$ 且 $C$ 由严格满态射生成.
再证此时 $E_{PF}$ 的每个对象 $X$ 都同构于 $C$ 的某个对象的直因子.
事实上, 已知 $X$ 是 $C$ 的对象在 $E$ 中的一个滤过归纳极限 $\varinjlim_iX_i$;
由 $E_{PF}$ 的定义, 对适当的 $i$, 典范同态 $X_i\rightarrow X$ 有左逆,
故 $X$ 确实是 $C$ 的对象 $X_i$ 的直因子.
这就证明了 (i) $\Rightarrow$ (ii) (此外, 不需对 $C$ 或 $E$ 作任何假设).
证明 (ii) $\Rightarrow$ (i). 由于 $C$ 等价于小范畴,
且 $E_{PF}$ 的每个对象都是 $C$ 的某个对象的直因子,
可见 $E_{PF}$ 也等价于小范畴.
故由 \hyperref[I.cor8.7.6]{8.7.6}, 函子
\hyperlink{sga4.exp.I.equation.8.7.6.1}{(8.7.6.1)} 为范畴等价,
再对包含 $E\hookrightarrow E_{PF}$ 应用 \hyperref[I.prop8.6.4]{8.6.4} b), 即得结论.

当 $C$ 的某个对象在 $E$ 中的每个直因子均属于 $C$ 时,
显然有 (ii) $\Leftrightarrow$ (ii bis).
另一方面, 由 \hyperref[I.cor8.7.6]{8.7.6}, (ii bis) 蕴含 (iii) 或分别蕴含 (iii bis).
最后, 由 \hyperref[I.prop8.7.5]{8.7.5}, (iii)
(从而更有 (iii bis)) 蕴含 (i). 证明完毕.

\begin{enonce*}[remark]{习题 8.7.8}\etiquette[8.7.8]{I.exer8.7.8}
(\emph{Karoubi 包络}). 设 $C$ 为 $𝒰$-范畴, $E=\Ind(C)$,
$E_{PF}\supset C$ 为 \hyperref[I.cor8.7.6]{8.7.6} 中定义的 $E$ 的满子范畴.
\begin{enumerate}
\renewcommand{\theenumi}{\alph{enumi}}
\renewcommand{\labelenumi}{\theenumi)}
\item 证明投影算子的像在 $E_{PF}$ 中可表,
且 $E_{PF}$ 的每个对象均为 $C$ 的某个对象的直因子.

\item 如果投影算子的像在范畴 $F$ 中可表, 就称它为 \emph{Karoubi 范畴}.
设
$$
\varphi:C\longrightarrow K
$$
为\pageoriginale{}全忠实函子, 使得 $K$ 为 Karoubi 范畴,
且 $K$ 的每个对象均为 $C$ 的像中某个对象的直因子.
证明 $\varphi$ 在从 $C$ 出发、取值于 Karoubi 范畴的函子中为 2-\emph{泛}的;
更确切地说, 对每个 Karoubi 范畴 $F$, 函子
$$
g\longmapsto g\circ \varphi :\Hhom(K,F)\longrightarrow \Hhom(K,F)
$$
为范畴等价. (利用每个函子都与投影算子的像交换这一事实.)
配备 $\varphi$ 的范畴 $K$ 由上述性质确定到等价
(该等价本身确定到唯一同构), 称为 $C$ 的 \emph{Karoubi 包络} $\widetilde C$.
(亦比较 IV \hyperref[I.cor7.5]{7.5}.)

\item 由 a) 和 b) 推出, $\Ind(C)_{PF}$ 配备由
$c:C\rightarrow\Ind(C)$ 诱导的函子 $C\rightarrow\Ind(C)_{PF}$ 后,
是 $C$ 的 Karoubi 包络.

\item 证明每个函子 $f:C\rightarrow C'$ 均以本质唯一的方式扩张为
Karoubi 包络的函子 $\widetilde f:\widetilde C\rightarrow\widetilde{C'}$,
且若如 c) 那样选取这些 Karoubi 包络, 则 $\widetilde f$ 是由
$\Ind(f):\Ind(C)\rightarrow\Ind(C')$ 诱导的函子.
证明 \hyperref[I.sec8.5.4]{8.5.4} %%
b) %%
的各条件还等价于: $\widetilde f$ 为范畴等价.
\end{enumerate}
\end{enonce*}

\begin{enonce*}[remark]{习题 8.7.9}\etiquette[8.7.9]{I.exer8.7.9}
设 $E$ 为 $𝒰$-范畴.
\begin{enumerate}
\renewcommand{\theenumi}{\alph{enumi}}
\renewcommand{\labelenumi}{\theenumi)}
\item 证明下述条件等价:

\begin{itemize}
\item[(i)] $E$ 等价于形如 $\Ind(C)$ 的范畴, 其中 $C$ 为 $𝒰$-范畴.

\item[(i bis)] 同\pageoriginale (i), 并进一步要求 $C$ 为 Karoubi 范畴
(\hyperref[I.exer8.7.8]{8.7.8} b)).

\item[(ii)] $E$ 的子范畴 $E_{PF}$ (\hyperref[I.cor8.7.6]{8.7.6})
由严格满态射生成, 且对 $E$ 的每个对象 $X$,
$(E_{PF})_{/X}$ 为本质小滤过范畴.
\end{itemize}

证明, 若 $C$ 为使 $E$ 等价于 $\Ind(C)$ 的 \emph{Karoubi} $𝒰$-范畴,
则 $C$ 等价于 $E_{PF}$ (此等价确定到唯一同构).
在由满足上述条件的范畴 $E$ 及它们之间\emph{与小滤过归纳极限交换}的函子
组成的 2-范畴, 与由 \emph{Karoubi} $𝒰$-范畴及它们之间任意函子组成的
2-范畴之间建立 2-等价.

\item 证明下述条件等价:

\begin{itemize}
\item[(i)] $E$ 等价于形如 $\Ind(C)$ 的范畴,
其中 $C$ 为有限归纳极限可表的 $𝒰$-范畴.

\item[(i bis)] 同 (i), 但进一步要求 $C$ 为 Karoubi 范畴.

\item[(ii)] $E$ 中的有限归纳极限可表, $E$ 的子范畴 $E_{PF}$ 由严格满态射生成,
且对 $E$ 的每个对象 $X$, 存在 $\Ob E_{PF/X}$ 的一个小子集 $I$,
使 $E_{PF/X}$ 的每个对象都以 $I$ 的某个对象为上界.
(使用 \hyperref[I.prop8.9.5]{8.9.5} b).)
\end{itemize}
\end{enumerate}
\end{enonce*}


% END INCLUDED SGA4_English_current_buildable_source_20260906/source/01/drafts/SGA4_I_LINES_4069_4270_English.texfrag


% BEGIN INCLUDED SGA4_English_current_buildable_source_20260906/source/01/drafts/SGA4_I_LINES_4271_4625_English.texfrag
